\documentclass[10pt,a4paper]{amsart}
\usepackage[margin=19mm]{geometry}
\usepackage{amsmath,amssymb,mathtools}
\usepackage[scr=rsfs,cal=euler]{mathalfa}
\usepackage{xfrac}
\usepackage[unicode,hidelinks,bookmarks=false]{hyperref}
\newcommand{\ie}{i.e.\ }
\newcommand{\cf}{cf.\ }
\newcommand{\resp}{respectively\ }
\newcommand{\Subsection}{\subsection}
\providecommand{\nobreakdash}{\nobreak\hbox{-}}
\setlength{\emergencystretch}{2em}
\multlinegap0pt
\usepackage{dsfont}
\usepackage{bm}
\usepackage{amssymb}
\usepackage{xcolor}
\usepackage{mathtools}
\usepackage[shortlabels]{enumitem}
\usepackage{tikz}
\newtheorem{definition}{Definition}
\newtheorem{theorem}{Theorem}
\newtheorem{lemma}{Lemma}
\def\M{{\mathcal{M}}}
\newcommand{\R}{\mathbb{R}}
\def\dd{{\, \rm d}}
\title{Nonlinear Stability and Instability of\\ Finite-Energy Solutions of the Compressible Euler-Riesz Equations with General Pressure Laws}
\author{Jose A. Carrillo, Samuel R. Charles, Gui-Qiang G. Chen, Difan Yuan}
\date{Source: arXiv:2607.26782v1 (CC BY 4.0)}
\begin{document}
\maketitle

\noindent\emph{Source and licence.} Nonlinear Stability and Instability of\\ Finite-Energy Solutions of the Compressible Euler-Riesz Equations with General Pressure Laws. Version arXiv:2607.26782v1 (CC BY 4.0), distributed by arXiv under the Creative Commons Attribution 4.0 International licence, http://creativecommons.org/licenses/by/4.0/, as recorded in the arXiv licence metadata for this exact version and shown on the abstract page https://arxiv.org/abs/2607.26782v1. Full licence notice: \texttt{licenses/arxiv-2607.26782-CC-BY-4.0.txt}.\par\smallskip

\noindent\textit{Editorial scope.} Counted unit: the article's theorem thm:stability (source0.tex lines 1000--1021). The article's complete original proof of it is delivered with it (source0.tex lines 3868--4206, one \texttt{proof} environment of 339 lines, which the article places away from the statement). No mathematics is written, repaired or re-derived: the statement and the proof are the article's own lines, in the article's own notation, and no retained line is substituted or re-flowed.  Retained uncounted as the source-local context the proof needs: lines 130--137; lines 231--234; lines 653--660; lines 930--944; lines 1147--1213; lines 3277--3279; lines 3688--3698; lines 3823--3832.  Nothing was omitted from any retained span, so the package carries no omission record.  No retained line contains a citation, so the wrapper carries no bibliography and no bibliography entry.  No retained line contains a figure environment, an image-inclusion command or drawing code, so no image is delivered and the package declares no asset dependency.  Editorial formatting only, in addition: an AMS \texttt{amsart} preamble that restates the article's own declarations and macros used by the retained lines, namely the environments \texttt{definition}, \texttt{theorem}, \texttt{lemma} on separate counters, together with the standard packages those lines need.  The retained environments print in this document as Definition 1, Theorem 1, Lemma 1 in order of appearance, whereas the article prints the same items with the numbering of its own full document.  The complete native source of the article as distributed in the arXiv e-print for this exact version is bundled unchanged as \texttt{sources/206293/source0.tex}.
\begin{align}\label{1.1}
\begin{cases}
\displaystyle
\partial_t \rho+\mbox{div} \mathcal{M}=0,\\
\displaystyle
\partial_t \mathcal{M}+\mbox{div}\Big(\frac{\mathcal{M}\otimes \mathcal{M}}{\rho}\Big)+\nabla p+\kappa\rho\nabla\mathcal{W}_\alpha=\boldsymbol{0},\\
\end{cases}
\end{align}

\begin{equation}\label{1.1-2}
(\rho, \M)|_{t=0}=(\rho_0, \M_0)( \boldsymbol{x})\longrightarrow (0, \boldsymbol{0})
\qquad\,\,\,\,\mbox{as $| \boldsymbol{x}|\to \infty$}.
\end{equation}

\begin{align}\label{station}
\begin{cases}
\frac{\dd}{\dd \rho}(\rho e(\rho))|_{\rho=\bar{\rho}}+\Phi_\alpha\ast\bar{\rho}= - \mu 
\quad & \mbox{for $\boldsymbol{x}\in \Gamma$},\\[1mm]
\Phi_\alpha\ast\bar{\rho}( \boldsymbol{x})\geq - \mu 
& \mbox{for $\boldsymbol{x}\in \R^n\setminus\Gamma$}.
\end{cases}
\end{align}

\begin{definition}\label{globalweak}
    A finite-energy solution $(\rho,\mathcal{M})$, as defined in Definition \ref{definition} below,
    is called an entropy solution if
    \[
 \partial_t \eta(\rho,\mathcal{M}) + \nabla \cdot q(\rho,\mathcal{M}) \leq -\kappa \mathcal{M}\cdot\nabla \mathcal{W}_\alpha
    \]
    in the sense of distributions for the mechanical entropy
    \[ \eta(\rho, \mathcal{M})
  = \frac{|\mathcal{M}|^{2}}{2\rho} + h(\rho),
  \qquad
  q(\rho,\mathcal{M})
  =\frac12 \mathcal{M}\frac{|\mathcal{M}|^{2}}{\rho^{2}}+\mathcal{M}h'(\rho),
\]
where $h(\rho) = \rho e(\rho)$.
\end{definition}

\begin{definition}\label{definition}
A measurable vector-valued function $(\rho,\M)(t, \boldsymbol{x})$ is a finite-energy weak solution 
of the Cauchy problem \eqref{1.1}--\eqref{1.1-2} if
$(\rho,\M)(t, \boldsymbol{x})$ satisfies the following conditions for all $T>0${\rm :}
\begin{enumerate}
\item[(a)] $\rho(t,\boldsymbol{x}) \geq 0$ {\it a.e.} and $(\M, \frac{\M}{\sqrt{\rho}})(t,\boldsymbol{x}) = \boldsymbol{0}$ {\it a.e.} on the vacuum states $\{ (t,\boldsymbol{x}) \, : \, \rho(t,\boldsymbol{x}) = 0 \}$,
\begin{align*}
\rho\in L^{\infty}(0,T; L^1(\mathbb{R}^n)), \quad \,\, 
\rho e(\rho)\in L^{\infty}(0,T; L^1(\mathbb{R}^n)), \quad\,\,
 \frac{\M}{\sqrt{\rho}}\in L^{\infty}(0,T; L^2(\mathbb{R}^n)).
    \end{align*}
 Moreover, for $\alpha \in (0,n)$,
    \[
\Phi_\alpha * \rho\in L^\infty(0,T;L^{\frac{2n}{\alpha}}(\mathbb{R}^n) ),\quad
\nabla\Phi_\alpha * \rho\in L^\infty(0,T;L^{\frac{2n}{\alpha + 2}}(\mathbb{R}^n)),
  \]
 with, for $\alpha \in [n-1,n)$,
 \[
 \rho \in L^\infty(0,T;C^{0,\beta}(\mathbb{R}^n)),
  \]
    for some $\beta \in (1 + \alpha - n,1)$; and, for $\alpha \in (-1,0]$,
    \[
 \Phi_\alpha * \rho\in L^\infty(0,T;L^\infty_{\rm loc}(\mathbb{R}^n) ),\quad
\nabla\Phi_\alpha * \rho\in L^\infty(0,T;L_{\rm loc}^{\frac{2n}{\alpha + 2}}(\mathbb{R}^n)).
 \]
\item[(b)] For {\it a.e.} $t>0$, the total energy is finite{\rm :}

\smallskip
\begin{itemize}
 \item For $\kappa = -1$,
\begin{equation*}
 \int_{\mathbb{R}^n} \Big( \frac{1}{2} \Big| \frac{\M}{\sqrt{\rho}}\Big|^2
 + \rho e(\rho) - \frac{1}{2} \rho (\Phi_\alpha * \rho) \Big)(t,\boldsymbol{x})
 \dd \boldsymbol{x} \leq E_0.
\end{equation*}
\item For $\kappa = 1$,
\begin{align*}
\begin{cases}
\displaystyle
\int_{\mathbb{R}^n} \Big( \frac{1}{2} \Big| \frac{\M}{\sqrt{\rho}}\Big|^2
+ \rho e(\rho) - \frac{1}{2} \rho (\Phi_\alpha * \rho) \Big)(t,\boldsymbol{x})
\dd \boldsymbol{x} \leq C(E_0,M), \\[4mm]
\displaystyle
\int_{\mathbb{R}^n} \Big( \frac{1}{2}\Big|\frac{\M}{\sqrt{\rho}}\Big|^2
+ \rho e(\rho) + \frac{1}{2} \rho (\Phi_\alpha * \rho) \Big)(t,\boldsymbol{x})
\dd \boldsymbol{x} \leq E_0.
\end{cases}
\end{align*}
\end{itemize}
\item[(c)] For any $\zeta \in C^1_0(\mathbb{R}_+ \times \mathbb{R}^n)$,
\begin{equation*}
\int_{\mathbb{R}_+} \int_{\mathbb{R}^n} \big(\rho \zeta_t
+ \M \cdot \nabla \zeta\big) \dd \boldsymbol{x} \dd t
+ \int_{\mathbb{R}^n} \rho_0(\boldsymbol{x})\,\zeta(0,\boldsymbol{x}) \dd \boldsymbol{x} = 0.
\end{equation*}
\item[(d)] For any $\boldsymbol{\psi} \in (C^1_0(\mathbb{R}_+ \times \mathbb{R}^n))^n$,
    \begin{align*}
 & \int_{\mathbb{R}_+} \int_{\mathbb{R}^n}
  \Big(\M \cdot \boldsymbol{\psi}_t
 + \frac{\M}{\sqrt{\rho}} \cdot \big( \frac{\M}{\sqrt{\rho}} \cdot \nabla \big)
 \boldsymbol{\psi} + p(\rho)\nabla \cdot \boldsymbol{\psi} \Big)(t,\boldsymbol{x})
 \dd \boldsymbol{x} \dd t
  + \int_{\mathbb{R}^n} \M_0(\boldsymbol{x}) \cdot \boldsymbol{\psi}(0,\boldsymbol{x}) \dd \boldsymbol{x} \\
  & = \int_{\mathbb{R}_+} \int_{\mathbb{R}^n} \kappa(\rho \nabla \mathcal{W}_\alpha \cdot \boldsymbol{\psi})(t,\boldsymbol{x}) \dd \boldsymbol{x} \dd t.
    \end{align*}
\end{enumerate}
\end{definition}

\begin{equation}\label{property}
G_\mu=-h'(\bar\rho)\quad\hbox{{\it a.e. }on } \Gamma, \qquad\,\, G_\mu\ge0\quad\hbox{{\it a.e. }on }\Gamma^{\rm c}.
\end{equation}

\begin{lemma}\label{relativenonlocal}
Under the assumptions of {\rm Theorem \ref{thm:stability}}, set
\[
f(t,\boldsymbol{x}):=\rho(t,\boldsymbol{x})-\bar\rho(\boldsymbol{x}).
\]
Then, for {\it a.e.} $t\in(0,T)$,
\begin{equation}\label{relativeestimate}
\int_{\mathbb R^n}\rho\,|\nabla\Phi_\alpha*f|^2\,\dd\boldsymbol{x}
 \le C\Big( \mathcal E_\mu(t) +\mathcal E_\mu(t)^{2/p_{\rm loc}} +\mathcal E_\mu(t)^{2/\gamma} +\mathcal E_\mu(t)^{3/2}+\mathcal E_\mu(t)^{3/p_0}  +\mathcal E_\mu(t)^{3/\gamma}  \Big).
\end{equation}
\end{lemma}

\begin{align}\label{testinequality}
&\int_0^T\!\!\int_{\mathbb R^n}
\Big(\partial_\tau\psi\,\eta_\mu(\rho,\mathcal M\mid\bar\rho)
+\nabla\psi\cdot q_\mu(\rho,\mathcal M\mid\bar\rho)\Big)\,\dd\boldsymbol{x}\dd\tau
+\int_{\mathbb R^n}\psi(0,\boldsymbol{x})\,\eta_\mu(\rho,\mathcal M\mid\bar\rho)(0,\boldsymbol{x})\,\dd\boldsymbol{x}
\nonumber\\
&\qquad\ge
\int_0^T\!\!\int_{\mathbb R^n}
\psi\,\mathcal M\cdot\nabla\Phi_\alpha*(\rho-\bar\rho)\,\dd\boldsymbol{x}\dd\tau.
\end{align}

\begin{theorem}[Weighted $L^2$--Stability Around Steady States]\label{thm:stability}
Let $\alpha\in(0,n-1)$, and let $(\bar\rho,\boldsymbol{0})$ be a compactly supported polytropic steady state characterized by \eqref{station}.
Let $(\rho,\mathcal M)$ be a global entropy solution of 
the attractive CEREs {\rm(}$\kappa=1${\rm)} \eqref{1.1} as in
{\rm Definition \ref{globalweak}}, and let
\[
\max \Big\{\frac{n + \alpha}{n},\frac{3n}{3n - 2(\alpha + 1)} \Big\}<\gamma\le2.
\]
Suppose that $\mathcal E_\mu(0)<\infty$ and
\[
(\rho^{\gamma + \theta}, \,\frac{|\mathcal{M}|^{3}}{\rho^{2}})
\in L^1(0,T;L^1(\mathbb{R}^n))
\qquad \mbox{for $\theta = \frac{\gamma - 1}{2}$}.
\]
Then, for every fixed $T>0$, there exists a constant
$C_T=C_T(\mathcal E_\mu(0),\bar{\rho},T)>0$ such that
\begin{equation}\label{finitetimerelative}
 \mathcal E_\mu(t)
 \le C_T\mathcal E_\mu(0)
 \qquad \hbox{for {\it a.e.} }t\in[0,T].
\end{equation}
\end{theorem}

\begin{proof}[Proof of {\rm Theorem \ref{thm:stability}}]
We divide the proof into six steps.

\smallskip
1. Fix a Lebesgue point $t\in(0,T)$.   Let $R_0>0$ be such that
$K\subset B_{R_0}$.  Choose $L>R_0+1$ and take
$0<\varepsilon<\min\{1,T-t\}$.  For such $\varepsilon$, define
\[
\vartheta_\varepsilon(\tau)=
\begin{cases}
1& \mbox{for $0\le\tau<t$},\\
1+(t-\tau)/\varepsilon& \mbox{for $t\le\tau<t+\varepsilon$},\\
0&\mbox{for $t+\varepsilon\le\tau\le T$},
\end{cases}
\]
and let $\zeta_L\in W^{1,\infty}_{\rm c}(\mathbb R^n)$ be chosen so that
\begin{align*}
&0\le \zeta_L\le1,
\qquad
\zeta_L=1\ \text{on }B_L,
\qquad
\zeta_L=0\ \text{on }B_{L+1}^{\rm c},\\[1mm]
&|\nabla\zeta_L|\le C\mathds{1}_{\{L<|\boldsymbol{x}|<L+1\}},
 \qquad
 \zeta_L(\boldsymbol{x}) \to 1^- 
 \quad\text{as }L\to\infty
 \quad\text{for every fixed }\boldsymbol{x}.
\end{align*}
Put $\psi_\varepsilon(\tau,\boldsymbol{x})=\vartheta_\varepsilon(\tau)\zeta_L(\boldsymbol{x})$ in \eqref{testinequality}.

\smallskip
2. For {\it a.e.} $\tau$, set $F(\tau):=\nabla\Phi_\alpha*(\rho(\tau)-\bar\rho)$.  Since $0\le\psi_\varepsilon\le1$,
\[
\begin{aligned}
\left|\int_{\mathbb R^n}\psi_\varepsilon\mathcal M\cdot F\,\dd\boldsymbol{x}\right|
&\le
\int_{\mathbb R^n}\frac{|\mathcal M|}{\sqrt\rho}\sqrt\rho |F|\,\dd\boldsymbol{x}\\
&\le
\frac12\int_{\mathbb R^n}\frac{|\mathcal M|^2}{\rho}\,\dd\boldsymbol{x}
+\frac12\int_{\mathbb R^n}\rho |F|^2\,\dd\boldsymbol{x}\\
&\le C\big(\mathcal E_\mu(\tau)+\mathcal E_\mu(\tau)^{\chi_*}\big),
\end{aligned}
\]
where Lemma \ref{relativenonlocal} has been used in the last inequality, 
and $\chi_*>1$.  Therefore, we have
\begin{equation}\label{sourceestimate}
\left|\int_0^T\!\!\int_{\mathbb R^n}
\psi_\varepsilon\mathcal M\cdot F\,\dd\boldsymbol{x}\dd\tau\right|
\le C\int_0^{t+\varepsilon}
\big(\mathcal E_\mu(\tau)+\mathcal E_\mu(\tau)^{\chi_*}\big)\,\dd\tau.
\end{equation}

\smallskip
3. Set $A_L:=\{L<|\boldsymbol{x}|<L+1\}$.  Notice that $\bar\rho=0$ on $A_L$. 
Moreover,
$
 |\nabla\zeta_L|\le C\mathds{1}_{A_L}.
$
On the same shell,
\[
 q_\mu(\rho,\mathcal M\mid\bar\rho)
=\frac12\mathcal M\frac{|\mathcal M|^2}{\rho^2}+h'(\rho)\mathcal M+G_\mu\mathcal M.
\]
Write $\mathcal M=\rho u$.  Then
\[
 \left|\frac12\mathcal M\frac{|\mathcal M|^2}{\rho^2}\right|
  =\frac12\rho |u|^3
  =\frac12\frac{|\mathcal M|^3}{\rho^2}.
\]
Since $h'(\rho)=a_0\gamma\rho^{\gamma-1}/(\gamma-1)$, $
 |h'(\rho)\mathcal M|
 \le C\rho^\gamma |u|.$
Using the Young inequality, we have
\[
\rho^\gamma|u|
 \le C\rho |u|^3+C\rho^{(3\gamma-1)/2}
 =C\frac{|\mathcal M|^3}{\rho^2}+C\rho^{\gamma+\theta}.
\]

Finally, since $K\subset B_{R_0}$ and $L>R_0+1$, 
we see that $|\boldsymbol{x}-\boldsymbol{y}|\ge 1$
for every
$\boldsymbol{x}\in A_L$ and every $\boldsymbol{y}\in K$.
Then the singularity of
$\Phi_\alpha$ is avoided uniformly in $L$.  Since $\bar\rho$ is compactly
supported, then
$\sup_{|\boldsymbol{x}|\ge R_0+1}
|\bar{\mathcal W}_\alpha(\boldsymbol{x})|<\infty.$
Consequently,
$\sup_{|\boldsymbol{x}|\ge R_0+1}|G_\mu(\boldsymbol{x})|<\infty.$
The corresponding bound is independent of $L$.  
Therefore, we obtain
\[
|G_\mu\mathcal M|
\le C\rho |u|
\le C\rho+C\rho |u|^3
=C\rho+C\frac{|\mathcal M|^3}{\rho^2}.
\]
Defining
\[
 \mathcal{F}(\tau,\boldsymbol{x}) :=\rho+\rho^{\gamma+\theta}+\frac{|\mathcal M|^3}{\rho^2},
\]
we have
\begin{equation}\label{fluxbound}
\left|\int_0^T\!\!\int_{\mathbb R^n}
\nabla\psi_\varepsilon\cdot q_\mu(\rho,\mathcal M\mid\bar\rho)\,\dd\boldsymbol{x}\dd\tau\right|
\le
C\int_0^T\!\!\int_{A_L} \mathcal{F}\,\dd\boldsymbol{x}\dd\tau.
\end{equation}
The right-hand side is finite because $\rho\in L^\infty(0,T;L^1(\R^n))$ for finite-energy solutions and, by assumptions,
$\rho^{\gamma+\theta}$ and $|\mathcal M|^3/\rho^2$ belong to $L^1(0,T;L^1(\R^n))$.

\smallskip
4. Since
\[
\partial_\tau\psi_\varepsilon =-\varepsilon^{-1}\zeta_L
 \,\,\,\,\hbox{on }(t,t+\varepsilon),
     \qquad
\partial_\tau\psi_\varepsilon=0
\,\,\,\,\hbox{elsewhere},
\]
\eqref{testinequality}, \eqref{sourceestimate}, and \eqref{fluxbound} imply
\begin{align}\label{epsiloninequality}
&\frac1\varepsilon\int_t^{t+\varepsilon}\!\!\int_{\mathbb R^n}
\zeta_L\eta_\mu(\rho,\mathcal M\mid\bar\rho)(\tau)\,\dd\boldsymbol{x}\dd\tau
\nonumber\\
&\quad\le
\int_{\mathbb R^n}\zeta_L(\boldsymbol{x})
\eta_\mu(\rho,\mathcal M\mid\bar\rho)(0,\boldsymbol{x})\,\dd\boldsymbol{x}
+C\int_0^{t+\varepsilon}\big(\mathcal E_\mu(\tau)+\mathcal E_\mu(\tau)^{\chi_*}\big)\,\dd\tau
\nonumber\\
&\qquad
+C\int_0^T\!\!\int_{A_L} \mathcal{F}\,\dd\boldsymbol{x}\dd\tau.
\end{align}
For fixed $L$, the Lebesgue differentiation theorem yields that, for {\it a.e.} such $t$,
\[
\lim_{\varepsilon\to0^+}
\frac1\varepsilon\int_t^{t+\varepsilon}\!\!\int_{\mathbb R^n}
\zeta_L\eta_\mu(\tau)\,\dd\boldsymbol{x}\dd\tau
=
\int_{\mathbb R^n}\zeta_L\eta_\mu(t)\,\dd\boldsymbol{x}.
\]
Letting $\varepsilon\to0^+$ in \eqref{epsiloninequality}, we obtain
\begin{align}\label{localL}
\int_{\mathbb R^n}\zeta_L\eta_\mu(t)\,\dd\boldsymbol{x}
&\le
\int_{\mathbb R^n}\zeta_L\eta_\mu(0)\,\dd\boldsymbol{x}
+C\int_0^t\big(\mathcal E_\mu(\tau)+\mathcal E_\mu(\tau)^{\chi_*}\big)\,\dd\tau
+C\mathcal R(L),
\end{align}
where
\[
 \mathcal R(L):=
 \int_0^T\!\!\int_{A_L} \mathcal{F}\,\dd\boldsymbol{x}\dd\tau.
\]

\smallskip
5. 
Since $ \mathcal{F}\in L^1((0,T)\times\mathbb R^n)$, we have
\[
 \mathcal R(L)\to0
 \qquad\text{as }L\to\infty .
\]
Moreover, $\zeta_L\nearrow1$ pointwise.  Hence, by the monotone convergence
theorem,
\[
\int_{\mathbb R^n}\zeta_L\eta_\mu(t)\,\dd\boldsymbol{x}\to\mathcal E_\mu(t),
  \qquad
\int_{\mathbb R^n}\zeta_L\eta_\mu(0)\,\dd\boldsymbol{x}\to\mathcal E_\mu(0).
\]
Letting $L\to\infty$ in \eqref{localL}, we obtain that, for {\it a.e.} $t\in(0,T)$,
\begin{equation}\label{ODEestimate} \mathcal E_\mu(t)\le
 \mathcal E_\mu(0)
+C\int_0^t\big(\mathcal E_\mu(\tau)+\mathcal E_\mu(\tau)^{\chi_*}\big)\,\dd\tau,
\end{equation}
where $C>0$ depends only on $n,\alpha,\gamma,a_0$, and the steady state.

\smallskip
6. We first prove that, for every fixed $T>0$, there exists a constant
$C_T=C_T(\mathcal E_\mu(0),\bar{\rho},T)>0$ such that
\begin{equation}\label{apriorrb}
 \operatorname*{ess\,sup}_{0\le t\le T}\mathcal E_\mu(t)\le C_T .
\end{equation}
Set
\[
 A(t):=\int_{\mathbb R^n}
 \big(\frac{|\mathcal M|^2}{2\rho}+h(\rho)\big)(t,\boldsymbol{x})\,
 \dd\boldsymbol{x}.
\]
Since $\Phi_\alpha\le0$ and $\bar\rho\ge0$, we see that
$\bar{\mathcal W}_\alpha\le0$.  Moreover, by \eqref{property},
\[
 G_\mu=-h'(\bar\rho)\le0
 \quad\hbox{\it a.e. on }\Gamma,
 \qquad
 0\le G_\mu=\mu+\bar{\mathcal W}_\alpha\le\mu
 \quad\hbox{\it a.e. on }\Gamma^{\rm c} .
\]
Hence, using $\bar\rho=0$ on $\Gamma^{\rm c}$, for {\it a.e.} $t\in(0,T)$,
\[
\begin{aligned}
 \int_{\mathbb R^n}G_\mu(\rho-\bar\rho)(t)\,\dd\boldsymbol{x}
 &=
 \int_{\Gamma}G_\mu\rho(t)\,\dd\boldsymbol{x}
 -\int_{\Gamma}G_\mu\bar\rho\,\dd\boldsymbol{x}
 +\int_{\Gamma^{\rm c}}G_\mu\rho(t)\,\dd\boldsymbol{x}  \\
 &\le
 -\int_{\Gamma}G_\mu\bar\rho\,\dd\boldsymbol{x}
 +\mu\int_{\Gamma^{\rm c}}\rho(t)\,\dd\boldsymbol{x}\le
 C_{\bar\rho}+\mu M .
\end{aligned}
\]
Therefore, we deduce
\[
\begin{aligned}
 \mathcal E_\mu(t)
 = A(t)-\int_{\mathbb R^n}h(\bar\rho)\,\dd\boldsymbol{x}
 +\int_{\mathbb R^n}G_\mu(\rho-\bar\rho)(t)\,\dd\boldsymbol{x}\le A(t)+C_{\bar\rho}+\mu M .
\end{aligned}
\]
By the energy estimate for finite-energy weak solutions,
\[
 \operatorname*{ess\,sup}_{0\le t\le T} A(t)
 \le C(A(0),M).
\]
Consequently, one obtains
\[
 \operatorname*{ess\,sup}_{0\le t\le T}\mathcal E_\mu(t)
 \le C(A(0),M,\bar\rho).
\]

It remains to express the right-hand side only in terms of
$\mathcal E_\mu(0)$ and $\bar\rho$.  Since
$G_\mu\ge0$ on $\Gamma^{\rm c}$ and $G_\mu=-h'(\bar\rho)$ on $\Gamma$,
\[
\begin{aligned}
 A(0)=\mathcal E_\mu(0)
   +\int_{\mathbb R^n}h(\bar\rho)\,\dd\boldsymbol{x}
   -\int_{\mathbb R^n}G_\mu(\rho_0-\bar\rho)\,\dd\boldsymbol{x}\le
 \mathcal E_\mu(0)+C_{\bar\rho}
 +\int_{\Gamma}h'(\bar\rho)\rho_0\,\dd\boldsymbol{x}.
\end{aligned}
\]
Since $\bar\rho\in L^\infty_{\rm c}(\mathbb R^n)$, $h'(\bar\rho)$ is bounded
and compactly supported.  By the Young inequality, for every $\delta>0$,
\[
 h'(\bar\rho)\rho_0
 \le \delta h(\rho_0)+C_{\delta,\bar\rho}
 \qquad\hbox{on }\Gamma .
\]
Therefore, we have
\[
 A(0)\le \mathcal E_\mu(0)+C_{\delta,\bar\rho}+\delta A(0).
\]
Choosing $\delta>0$ sufficiently small and absorbing the last term into the
left-hand side give
\[
 A(0)\le C_{\bar\rho}\bigl(1+\mathcal E_\mu(0)\bigr).
\]

We also control the mass $M$ by $\mathcal E_\mu(0)$.  Since
$\bar\rho$ is compactly supported and
$\bar{\mathcal W}_\alpha(\boldsymbol{x})\to0$ as
$|\boldsymbol{x}|\to\infty$, there exists $R>0$, depending only on
$\bar\rho$ and $\mu$, such that
\[
 K\subset B_R
 \qquad\hbox{and}\qquad
 G_\mu(\boldsymbol{x})
 =\mu+\bar{\mathcal W}_\alpha(\boldsymbol{x})
 \ge \frac{\mu}{2}
 \,\,\,\,\hbox{for }|\boldsymbol{x}|\ge R .
\]
In particular, $B_R^{\rm c}\subset\Gamma^{\rm c}$.  
On $\Gamma$, by the convexity of
$h$,
\[
 h_\mu(\rho_0\mid\bar\rho)
 =h(\rho_0)-h(\bar\rho)-h'(\bar\rho)(\rho_0-\bar\rho)\ge0,
\]
while, on $\Gamma^{\rm c}$,
\[
 h_\mu(\rho_0\mid0)=h(\rho_0)+G_\mu\rho_0\ge G_\mu\rho_0 .
\]
Since the kinetic part is nonnegative, we obtain
\[
 \mathcal E_\mu(0)
 \ge
 \int_{B_R^{\rm c}}G_\mu\rho_0\,\dd\boldsymbol{x}
 \ge
 \frac{\mu}{2}\int_{B_R^{\rm c}}\rho_0\,\dd\boldsymbol{x}.
\]
Thus, we infer that
\[
 \int_{B_R^{\rm c}}\rho_0\,\dd\boldsymbol{x}
 \le C_\mu\mathcal E_\mu(0).
\]
Moreover, since $B_R$ has finite measure,
\[
 \int_{B_R}\rho_0\,\dd\boldsymbol{x}
 \le C_R\Big(1+\int_{B_R}h(\rho_0)\,\dd\boldsymbol{x}\Big)
 \le C_{\bar\rho}\bigl(1+A(0)\bigr)
 \le C_{\bar\rho}\bigl(1+\mathcal E_\mu(0)\bigr).
\]
Therefore, we deduce
\[
 M=\int_{\mathbb R^n}\rho_0\,\dd\boldsymbol{x}
 \le C_{\bar\rho}\bigl(1+\mathcal E_\mu(0)\bigr).
\]
Combining the estimates for $M$ and $A(0)$,  we have
\[
 \operatorname*{ess\,sup}_{0\le t\le T}\mathcal E_\mu(t)
 \le C(\mathcal E_\mu(0),\bar\rho).
\]
This proves \eqref{apriorrb}.


Since $\mathcal E_\mu(t)\le C$ for {\it a.e.} $t\in[0,T]$ and $\chi_*>1$, we have
\[
 \mathcal E_\mu(t)^{\chi_*}
 \le C^{\chi_*-1}\mathcal E_\mu(t)
 \qquad\hbox{for {\it a.e.} }t\in[0,T].
\]
Substituting this into \eqref{ODEestimate} yields
\[
 \mathcal E_\mu(t)
 \le \mathcal E_\mu(0)
 +C\bigl(1+C^{\chi_*-1}\bigr)\int_0^t\mathcal E_\mu(\tau)\,\dd\tau\qquad 
\hbox{for {\it a.e.} }t\in[0,T].
\]
By the Grönwall inequality,
\[
 \mathcal E_\mu(t)
 \le \mathcal E_\mu(0)\exp\{C(1+C^{\chi_*-1})t\}
 \le C_T\mathcal E_\mu(0)
 \qquad \hbox{for {\it a.e.} }t\in[0,T],
\]
where $C_T=C_T(\mathcal E_\mu(0), \bar{\rho}, T)$.  This proves \eqref{finitetimerelative}.
\end{proof}

\end{document}
